\documentclass[3p,times,twocolumn]{elsarticle}

% Elsevier's class requires the title, authors, abstract and keywords inside a frontmatter block,
% which pandoc's default template does not emit; everything else here is the usual pandoc preamble.
\usepackage{amsmath,amssymb}
\usepackage{graphicx}
\usepackage{longtable,booktabs,array}
\usepackage{calc}
\usepackage{float}
\usepackage[normalem]{ulem}
\usepackage{xcolor}
\usepackage[hidelinks]{hyperref}
\providecommand{\tightlist}{\setlength{\itemsep}{0pt}\setlength{\parskip}{0pt}}
\providecommand{\pandocbounded}[1]{#1}
\providecommand{\passthrough}[1]{#1}
\setlength{\emergencystretch}{3em}
\providecommand{\passthrough}[1]{#1}
\providecommand{\tightlist}{\setlength{\itemsep}{0pt}\setlength{\parskip}{0pt}}
\providecommand{\pandocbounded}[1]{#1}

% wide floats in two-column mode: a single-column figure of a full-width panel is unreadable
\makeatletter
\renewenvironment{figure}[1][]{\@float{figure}[#1]}{\end@float}
\makeatother

% symbols used in the Markdown source that the default text font lacks
\usepackage{newunicodechar}
\newunicodechar{≥}{\ensuremath{\geq}}
\newunicodechar{≤}{\ensuremath{\leq}}
\newunicodechar{∞}{\ensuremath{\infty}}
\newunicodechar{−}{\ensuremath{-}}

\journal{Engineering Applications of Artificial Intelligence}

\begin{document}
\begin{frontmatter}

\title{Reliability-gated mixture of experts with conformal calibration:
an uncertainty-aware digital twin for orbital debris, atmospheric
re-entry and satellite health}

\author{{[}Author name(s){]}\textsuperscript{a}}
\author{\textsuperscript{a}{[}Affiliation, postal address, country{]}}
\author{Corresponding author: {[}name, e-mail{]}}

\begin{abstract}
\emph{Contribution in artificial intelligence.} A reliability-gated
mixture of experts in which a learned gate weights a physics-based prior
against recent tabular neural models, with distribution-free conformal
calibration of every predictive interval and probability.
\emph{Application in engineering.} Forecasting the low-Earth-orbit
debris environment: residual orbital lifetime, collision risk and
spacecraft health.

Operational debris tools rely on physics propagators whose uncertainty
is uncalibrated, and prior machine-learning lifetime studies used a few
hundred to 1,800 objects, single-epoch predictions and only objects that
had already decayed. The proposed digital twin holds one state per
object and couples three learning modules through four explicit
mechanisms. Evaluated on 64,516 catalogued objects across two
forward-in-time test periods, a European Space Agency
collision-avoidance benchmark and public spacecraft telemetry, the
tabular foundation models reduced re-entry pinball loss 13 \% below the
strongest gradient-boosted baseline and 44 \% below the physics expert
for 2016--2021, yet were statistically tied with it for 2011--2015.
Conformal intervals covered 72.6 \%--90.6 \% of decayed objects against
12.2 \%--17.1 \% for the operational ±20 \% rule, and conditioning each
area-to-mass ratio on catalogue history reduced 72-day over-prediction
from 9.5-fold to 1.4-fold. Collision-risk forecasts improved 37 \% over
persistence, yet none beat the latest-risk baseline on the official
metric, and zero-shot telemetry detection did not beat simple detectors.
Calibrated learned fusion therefore improves debris forecasting wherever
the full predictive distribution matters.

\textbf{Keywords:} Mixture of experts; Conformal prediction; Tabular
foundation models; In-context learning; Uncertainty quantification;
Digital twin; Space debris; Atmospheric re-entry prediction
\end{abstract}


\end{frontmatter}

\section{Introduction}\label{sec:intro}

More than 30,000 tracked objects share low Earth orbit with the
satellites that operators must keep healthy, steer clear of debris and
eventually bring down. Three questions govern the resulting debris
environment. How long does each object remain in orbit? How likely is a
given encounter to become a collision? And which objects are still under
control, and therefore not decaying freely at all? In operational
practice each question is answered by a separate tool, and the
uncertainty those tools report is rarely calibrated.

The three questions are not independent. A spacecraft that loses the
ability to manoeuvre stops station-keeping and begins to decay;
fragments of one breakup share an origin, and therefore a range of
area-to-mass ratios that shapes how long they stay in orbit; and how
long fragments remain aloft determines the population that conjunction
assessment must screen. A digital twin of the tracked population, with
one state per object that every function reads and updates, is the
natural way to make these links explicit \citep{liu2024digital}.

Building such a twin on public data meets a hard constraint: no public
dataset follows the same object through telemetry, conjunction messages
and re-entry. Telemetry benchmarks are anonymised
\citep{kotowski2026esaadb, hundman2018}, conjunction messages are
anonymised \citep{uriot2022spacecraft}, and catalogue data carry no
telemetry. A credible twin therefore cannot be a single model trained
over unpaired data. It must be a set of modules, each validated on its
own real data, joined by couplings that are defined explicitly and
tested wherever a real test exists.

This paper presents such a system. Its contribution in artificial
intelligence is a \textbf{reliability-gated mixture of experts with
distribution-free calibration}: a learned gate weights a physics-based
prior against recent tabular neural models, and every interval and
probability the system emits is calibrated by conformal prediction and
isotonic regression. Its application in engineering is
\textbf{forecasting the low-Earth-orbit debris environment}, across
residual orbital lifetime, collision risk and spacecraft health.

The contributions are:

\begin{enumerate}
\def\labelenumi{\arabic{enumi}.}
\tightlist
\item
  A reliability-gated mixture-of-experts architecture for censored
  regression on an operational catalogue, evaluated against eight
  alternative fusion arms including the physics prior alone, each
  learned expert alone, an equal-weight average, a learned gate without
  reliability features, and a gradient-boosted baseline.
\item
  A licence-clean, leakage-controlled benchmark of 64,516 catalogued
  objects with administrative censoring and two independent
  forward-in-time test periods.
\item
  Four explicit cross-module couplings, three of which are tested on
  real data, with their measured contribution and their statistical
  significance reported.
\item
  Evidence that conformal prediction's marginal guarantee degrades under
  temporal drift, and that a catalogue covariate can leak the outcome
  through its availability rather than its value.
\item
  An honest account of where learning helps and where it does not,
  including a central negative result: the learned gate never beat the
  best single expert.
\end{enumerate}

\section{Related work and the research gap}\label{sec:related}

\subsection{Physics-based lifetime and re-entry
prediction}\label{physics-based-lifetime-and-re-entry-prediction}

Semi-analytic decay theory
\citep{kinghele1987satellite, frey2019extension} and numerical
propagators with empirical atmospheres underpin operational tools such
as the European Space Agency's debris-risk suite \citep{braun2020drama}.
Assessments of the inter-agency re-entry campaigns report mean errors of
about 10 \% of the remaining lifetime, and a ±20 \% window around the
predicted epoch that covers roughly 90 \% of true epochs in the final
days before re-entry \citep{pardini2018assessing}; earlier assessments
of propagator accuracy reach similar conclusions
\citep{pardini2004accuracy, operti2017assessing}. Thermospheric density
is the dominant error source, and both its empirical modelling
\citep{emmert2021nrlmsis, licata2022msis} and its uncertainty
\citep{licata2022uncertainty, pardini2012thermospheric} are active
subjects. Where orbital histories are available, the ballistic
coefficient is fitted from tracking data
\citep{gondelach2017ballistic, sang2013estimation, pardini2008impact},
which is the practice our state-estimation step approximates from two
catalogue orbits.

\subsection{Learning-based lifetime and re-entry
prediction}\label{learning-based-lifetime-and-re-entry-prediction}

Machine-learning studies predict the re-entry epoch or the orbital
lifetime from orbital histories. Table \ref{tbl:gap} summarises those we
could read in full. Salmaso et al.~train a sequence model on 417
Space-Track bodies and test on 15 campaign objects
\citep{salmaso2023machine}; Lotti et al.~fit a physics-informed
convolutional network to 1,783 payloads and report that generalisation
is ``not yet fully satisfactory'' \citep{lotti2025exploring}; Jung et
al.~use a recurrent network on altitude histories
\citep{jung2021recurrent}, and a physics-based backtest reports
calibrated lifetime distributions for 1,283 objects but calibrates on
the same backtest set \citep{shambaugh2026validation}. Related learning
work addresses density correction \citep{licata2022machine},
ballistic-coefficient estimation \citep{andaristiyan2025ballistic} and
lifetime regression on orbital elements
\citep{li2023backpropagation, peng2020machine}.

\begin{table*}[t]
\centering
\footnotesize
\caption{\textbf{Prior machine-learning work on orbital lifetime,
against this study.} Cells are taken from our reading of each source;
where a paper could not be read in full, the cell says so rather than
guessing.}\label{tbl:gap}
\begin{tabular}{@{}
>{\raggedright\arraybackslash}p{(\linewidth - 10\tabcolsep) * \real{0.1667}}
  >{\raggedright\arraybackslash}p{(\linewidth - 10\tabcolsep) * \real{0.1667}}
  >{\raggedright\arraybackslash}p{(\linewidth - 10\tabcolsep) * \real{0.1667}}
  >{\raggedright\arraybackslash}p{(\linewidth - 10\tabcolsep) * \real{0.1667}}
  >{\raggedright\arraybackslash}p{(\linewidth - 10\tabcolsep) * \real{0.1667}}
  >{\raggedright\arraybackslash}p{(\linewidth - 10\tabcolsep) * \real{0.1667}}@{}}
\toprule\noalign{}
\begin{minipage}[b]{\linewidth}\raggedright
Study
\end{minipage} & \begin{minipage}[b]{\linewidth}\raggedright
Scale
\end{minipage} & \begin{minipage}[b]{\linewidth}\raggedright
Censoring
\end{minipage} & \begin{minipage}[b]{\linewidth}\raggedright
Output
\end{minipage} & \begin{minipage}[b]{\linewidth}\raggedright
Evaluation split
\end{minipage} & \begin{minipage}[b]{\linewidth}\raggedright
Cross-module couplings
\end{minipage} \\
\midrule\noalign{}

\bottomrule\noalign{}

Salmaso et al. \citep{salmaso2023machine} & 417 bodies; test 15 campaign
objects & decayed only & point epoch & interleaved in time & none \\
Lotti et al. \citep{lotti2025exploring} & 1,783 after filtering; 1,435
train & decayed only & point lifetime & random, stratified & none \\
Jung et al. \citep{jung2021recurrent} & not readable (closed access) &
not readable & point epoch & not readable & none \\
Shambaugh \citep{shambaugh2026validation} & 1,283 objects & decayed only
& distribution & backtest, calibrated on itself & none \\
This work & 64,516 objects, two periods & censoring-aware & seven
quantiles + intervals & forward in time, two origins & three tested \\
\end{tabular}
\end{table*}

Three properties are shared by every study we could read. The datasets
are one to two orders of magnitude smaller than an operational
catalogue; the output is a point epoch rather than a distribution; and
the training population consists of objects that have already decayed,
which discards the censored majority. Our Kaplan--Meier description
(Section \ref{sec:ablations}) shows that this majority is 46.7 \% of one
test cohort and 55.3 \% of the other. None of these studies couples
lifetime prediction to conjunction assessment or to spacecraft health.

\subsection{Collision risk and the debris
environment}\label{collision-risk-and-the-debris-environment}

The Kelvins collision-avoidance challenge released conjunction data
messages with the final risk of each event
\citep{uriot2021cacdataset, uriot2022spacecraft}. The latest-risk
baseline scored 0.694 on the official metric and only 12 of 97 teams
beat it; the winning entry scored 0.556 \citep{uriot2022spacecraft}.
Later entries report scores on the same split
\citep{tulczyjew2021predicting}, and a pre-registered re-evaluation
found that no model beats the latest-risk forecast with confidence
\citep{vazquez2026latest}. Other learning approaches to conjunction
screening and risk include surrogate models of the encounter geometry
\citep{acciarini2021kessler, acciarini2024improving}, sequence models
over message streams \citep{chu2024conjunction, chu2026conjunction}, and
operational studies of alert handling
\citep{stevenson2022towards, lidtke2019optimising}. The debris
environment itself is modelled from breakup statistics
\citep{johnson2001nasa, kessler1978collision}, and breakup families have
been studied for origin attribution
\citep{wu2023investigation, celletti2021reconnecting}; the Kelvins
``space debris: the origin'' data are simulated
\citep{efthymiopoulos2021origin, esa2021kelvinsorigin}. We found no
study that uses breakup-family membership as a predictor of fragment
lifetime.

\subsection{Spacecraft telemetry anomaly
detection}\label{spacecraft-telemetry-anomaly-detection}

The public spacecraft benchmarks are the Soil Moisture Active Passive
and Mars Science Laboratory channels released with telemanom
\citep{hundman2018}, the European Space Agency anomaly-detection
benchmark \citep{kotowski2026esaadb, esaadb_dataset} and the OPS-SAT set
\citep{ruszczak2025opssat}. Metric design is a known hazard in this
literature: point adjustment can make random detectors appear strong
\citep{kim2022rigorous, wu2021flawed}, and range-based metrics were
introduced to address it \citep{tatbul2018precision}. Broad surveys of
time-series anomaly detection make the same point about benchmark
artefacts
\citep{schmidl2022anomaly, paparrizos2022tsbuad, sehili2023flawed}. We
therefore report event-wise and point-wise metrics without adjustment,
alongside label-free references.

\subsection{Methods from artificial
intelligence}\label{methods-from-artificial-intelligence}

Tabular prediction has moved from gradient-boosted trees
\citep{chen2016xgboost, ke2017lightgbm, prokhorenkova2018catboost},
which remained competitive on tabular benchmarks
\citep{grinsztajn2022tree}, to parameter-efficient deep ensembles
\citep{gorishniy2025tabm} and to in-context foundation models that
condition on the training set at inference time rather than fitting
weights to it
\citep{hollmann2025accurate, qu2025tabicl, qu2026tabiclv2, nvidia2026kumo}.
For time series, pretrained forecasters are used zero-shot
\citep{ansari2024chronos, ansari2025chronos2}. These models are normally
assessed on curated benchmark suites; we are not aware of an evaluation
on an operational scientific catalogue with administrative censoring and
forward-in-time drift.

Uncertainty is handled here by conformal prediction, which gives
finite-sample, distribution-free coverage under exchangeability
\citep{vovk2005algorithmic, angelopoulos2023gentle}. Conformalised
quantile regression adapts interval width to the input
\citep{romano2019cqr}, Mondrian variants condition coverage on groups
\citep{vovk2005algorithmic}, and the behaviour of the guarantee under
distribution shift has been studied directly
\citep{barber2023, candes2023conformalized, gibbs2021aci}. Conformal
methods have been applied to space-surveillance tracking
\citep{laxhammar2010conformal, laxhammar2015inductive}. Averaging
quantile functions is a well-studied combination rule
\citep{lichtendahl2013average}, and deep ensembles are a standard
uncertainty baseline \citep{lakshminarayanan2017ensembles}. Censored
targets connect the problem to survival analysis
\citep{harrell1982evaluating, barnwal2022survival}. Digital twins for
space systems have been proposed mostly as architectures and reviews
\citep{liu2024digital, guimaraes2023taxonomy}, or built for single
functions such as orbit prediction \citep{xu2025digital} and
physics-informed monitoring \citep{em2026physics}.

\subsection{The gap this paper
addresses}\label{the-gap-this-paper-addresses}

Three gaps follow. First, no catalogue-scale, censoring-aware,
forward-in-time evaluation of learned lifetime prediction with
calibrated intervals exists. Second, recent tabular and time-series
foundation models have not been tested on operational space-surveillance
data, where the covariate distribution drifts with the solar cycle.
Third, and most relevant to an integrated system, the couplings between
health, debris and re-entry are asserted in digital-twin architectures
but, as far as we can establish, never measured. This paper addresses
all three, and reports the couplings that do not pay off as well as
those that do.

\section{Objectives and problem formulation}\label{sec:objectives}

\subsection{Objectives}\label{objectives}

\begin{enumerate}
\def\labelenumi{\arabic{enumi}.}
\tightlist
\item
  To build a shared-state digital twin in which orbital debris analysis
  is the primary function, and lifetime prediction and spacecraft health
  are sub-functions that feed it.
\item
  To design a fusion architecture that combines a physics prior with
  recent tabular neural models and emits calibrated predictive
  distributions rather than point estimates.
\item
  To measure what each learned component and each cross-module coupling
  contributes, against the operational practice it would replace, under
  forward-in-time evaluation.
\item
  To report honestly where learning does not help.
\end{enumerate}

\subsection{Problem formulation}\label{problem-formulation}

\textbf{Residual lifetime.} For object \(i\) with state \(s_i(t_0)\) at
prediction time \(t_0\), let \(T_i\) be the time from \(t_0\) to natural
atmospheric re-entry. \(T_i\) is right-censored: an object still in
orbit at the catalogue snapshot, or whose phase ended by a controlled
de-orbit, landing or docking, contributes only the information
\(T_i > c_i\) for its censoring time \(c_i\). With a five-year horizon
the learning target is

\begin{equation}\protect\phantomsection\label{eq:target}{z_i = \log_{10} \min(T_i,\, 5\ \text{yr}), \qquad \delta_i = \mathbb{1}[T_i \leq c_i],}\end{equation}

and the estimand is the conditional distribution of \(z_i\) given
\(s_i(t_0)\), represented by the quantiles \(q_\tau(i)\) for
\(\tau \in \{0.05, 0.1, 0.25, 0.5, 0.75, 0.9, 0.95\}\) and by the
horizon probabilities \(P(T_i \leq H)\). Models are scored only on
objects whose five-year outcome is known at the snapshot, so the
five-year cap is a genuine point mass of the target rather than a
censoring artefact.

\textbf{Collision risk.} For encounter \(e\) with the sequence of
conjunction data messages \(C_e\) released at least two days before
closest approach, the target is the risk \(r_e\) of the final message,
where \(r_e = \log_{10} P_c\) is the base-ten logarithm of the collision
probability, floored at \(-30\). The estimand is the distribution of
\(r_e\) given \(C_e\), and the derived decision quantity is
\(P(r_e \geq -6)\).

\textbf{Health.} For channel \(c\) of a spacecraft with nominal history
\(\mathcal{C}_c\), a forecaster produces a non-conformity score
\(s_c(t)\), and the estimand is a calibrated alarm: a decision rule
whose rate of alarm on nominal data is controlled.

\textbf{What integration must mean.} We require each coupling to be a
function between module outputs and the shared state, with its inputs
restricted to information available at the prediction time. A coupling
counts as validated only if removing or replacing it changes a measured
quantity on real data by an amount whose confidence interval we report.
Architectural connections that cannot be tested on public data are
labelled as demonstrations and excluded from every performance claim.

\section{Architecture}\label{sec:arch}

\subsection{Shared state}\label{shared-state}

The system keeps one state per tracked object,

\begin{equation}\protect\phantomsection\label{eq:state}{s_i(t) = \bigl(o_i(t),\, \pi_i,\, h_i(t),\, F_i(t),\, e(t),\, C_i(t)\bigr),}\end{equation}

holding the latest orbit \(o_i\) (perigee, apogee, inclination), design
properties \(\pi_i\) (mass, dimensions, shape and a mass-uncertainty
flag), health evidence \(h_i\) (catalogue operational status, and an
alarm state where telemetry exists), the debris family \(F_i\) with
sibling statistics restricted to events before \(t\), space weather
\(e\) from observed days only, and a conjunction summary \(C_i\) where
messages exist. Three modules read and write this state, and every
output written to it carries its model, its validation status and its
own uncertainty. Figure \ref{fig:arch} shows the modules and Figure
\ref{fig:integration} the couplings.

\begin{figure*}
\centering
\includegraphics[width=0.92\linewidth,height=\textheight,keepaspectratio,alt={System architecture. Three modules read and update one state per object. Solid arrows are couplings tested on real data; the dashed coupling has no public test data.}]{artwork/Figure_4.png}
\caption{\textbf{System architecture.} Three modules read and update one
state per object. Solid arrows are couplings tested on real data; the
dashed coupling has no public test data.}\label{fig:arch}
\end{figure*}

\begin{figure*}
\centering
\includegraphics[width=1\linewidth,height=\textheight,keepaspectratio,alt={Integration. The four couplings as a data flow, with the variable that crosses each boundary and how each is evaluated. Operational status reaches the re-entry module through the shared state, which is how it is implemented.}]{artwork_eaai/eaai_integration.png}
\caption{\textbf{Integration.} The four couplings as a data flow, with
the variable that crosses each boundary and how each is evaluated.
Operational status reaches the re-entry module through the shared state,
which is how it is implemented.}\label{fig:integration}
\end{figure*}

\subsection{The artificial-intelligence
contribution}\label{the-artificial-intelligence-contribution}

The contribution is a two-expert mixture whose mixing weight is learned
from descriptors of \emph{one expert's own reliability}, and whose
output is calibrated without distributional assumptions. For object
\(i\),

\begin{equation}\protect\phantomsection\label{eq:gate}{q_\tau(i) = g(r_i)\, q^{\text{phys}}_\tau(i) + \bigl(1 - g(r_i)\bigr) q^{\text{ml}}_\tau(i), \qquad g(r) = \sigma(w^\top r + b),}\end{equation}

where \(q^{\text{phys}}_\tau\) are the quantiles of a Monte-Carlo
physics prior, \(q^{\text{ml}}_\tau\) those of a learned expert, and
\(r_i\) collects reliability descriptors of the physics prior: the
provenance of its area-to-mass ratio, the spread of its own Monte-Carlo
distribution, eccentricity, a perigee band, the object class, and the
probability it places at the five-year cap. Because a convex combination
of monotone quantile functions is monotone, Eq. \ref{eq:gate} returns a
valid predictive distribution for any \(g\); no sorting or post-hoc
repair is required. The parameters \((w,b)\) minimise the mean pinball
loss on out-of-fold predictions of the training period with an \(L_2\)
penalty, so the gate never sees a target that its experts were fitted
on.

Two design choices distinguish this from a standard stacked ensemble.
First, the gate input describes the \emph{prior's} trustworthiness, not
the input in general, which makes the learned weights interpretable: a
positive coefficient means ``trust the physics more in this regime''.
Second, mixing happens in quantile space rather than in expectation, so
the fused object remains a distribution and can be conformalised
directly.

\subsection{Uncertainty layer}\label{uncertainty-layer}

Two distinct calibrations are applied, and the distinction matters.
\textbf{Mondrian conformal prediction} converts the fused quantiles into
intervals with finite-sample coverage: on a calibration split disjoint
from fitting and gating, the score
\(S = \max(q_{\text{lo}} - z,\; z - q_{\text{hi}})\) is computed per
object, and within group \(k\) the correction \(\hat{Q}_k\) is the
\(\lceil (1-\alpha)(n_k+1) \rceil\)-th smallest score, giving the
interval \([q_{\text{lo}} - \hat{Q}_k,\; q_{\text{hi}} + \hat{Q}_k]\).
Groups are object class crossed with perigee band, falling back to class
and then to the pooled cohort when a group is too small.
\textbf{Isotonic regression} is a separate, monotone recalibration of
the horizon probabilities \(P(T \leq H)\) read off the same fused
quantile function. Conformal prediction controls interval coverage and
makes no assumption about the score distribution; isotonic regression
corrects the probability scale and assumes only monotonicity. They are
applied to different outputs and neither substitutes for the other.
Because isotonic maps are piecewise constant and therefore tie scores,
every ranking metric in this paper is computed on the raw probabilities
and every calibration metric on the mapped ones.

\section{Debris module: collision risk and cloud
population}\label{sec:debris}

\subsection{Family evidence}\label{family-evidence}

For object \(i\) at time \(t\) the module computes \(\varphi_F(i,t)\):
the number of siblings catalogued by \(t\) in the same launch family
and, for fragments, the same breakup family; the number and fraction of
those siblings that decayed naturally before \(t\); and their mean
log-lifetime. Only events strictly before \(t\) enter, so the feature is
causal by construction. This is the quantity that coupling C1 passes to
the re-entry module.

\subsection{Collision risk}\label{collision-risk}

Inputs are the conjunction data messages of an encounter released at
least two days before closest approach, reduced to one feature row per
event: the latest risk; the minimum, maximum, mean and first risk; the
change and the three-point slope; the last step; the logarithmic ratio
of miss distance and of the covariance scales between the first and the
last message; the message count; and the span. The target is the risk of
the final message.

A \textbf{persistence expert} predicts the latest risk plus the training
distribution of (final \(-\) latest) changes within the same latest-risk
bin, which is an honest persistence forecast with uncertainty rather
than a bare point. The same learned experts as the re-entry module are
applied to the event features, and the two are fused by Eq.
\ref{eq:gate}. Learner selection, the gate, the conformal calibration,
the isotonic map and the decision threshold are all driven by nested
out-of-fold predictions over five event-hash folds, so no test message
informs any choice. Quantiles within \(10^{-4}\) of the risk floor are
snapped to it, because most final risks lie exactly at the floor and the
spacing of single-precision floats there is about \(2 \times 10^{-6}\).

Because the official metric scores a \emph{point} forecast, a decision
rule is also required: predict the latest risk when the calibrated
\(P(r \geq -6)\) exceeds a threshold chosen on out-of-fold training
events, and \(-6.001\) otherwise.

\subsection{Cloud population}\label{cloud-population}

Coupling C2 turns per-object survival into a debris-environment
quantity. For a breakup cloud the expected number of catalogued
fragments still in orbit at time \(t\) is

\begin{equation}\protect\phantomsection\label{eq:cloud}{\mathbb{E}[N(t)] = \sum_{i:\, t_{0,i} \leq t} S_i\bigl(t - t_{0,i}\bigr),}\end{equation}

where \(S_i\) is the survival function of fragment \(i\) read from the
fused predictive distribution. It is read two ways: directly from the
quantile function, and after isotonic maps of \(P(T \leq a)\) fitted on
calibration fragments at 25 ages and interpolated in log-age. Reading it
two ways is necessary because the two readings disagree, as Section
\ref{sec:results} shows.

\section{Re-entry module: physics prior and learned
expert}\label{sec:reentry}

Figure \ref{fig:method} shows the whole pipeline, including which cohort
each component is fitted on.

\begin{figure*}
\centering
\includegraphics[width=1\linewidth,height=\textheight,keepaspectratio,alt={Experimental design of the re-entry module. Every component is fitted on data strictly earlier than the cohort it is applied to, and the test cohort enters only at the final step.}]{artwork_eaai/eaai_method.png}
\caption{\textbf{Experimental design of the re-entry module.} Every
component is fitted on data strictly earlier than the cohort it is
applied to, and the test cohort enters only at the final
step.}\label{fig:method}
\end{figure*}

\subsection{Physics prior}\label{physics-prior}

The prior integrates the orbit-averaged drag equations for the
semi-major axis \(a\) and eccentricity \(e\) with a thermospheric
density from NRLMSIS 2.1 \citep{emmert2021nrlmsis} evaluated at perigee,
and a ballistic coefficient \(B = C_D A/m\). For \(e \leq 0.05\) it uses
King-Hele's Bessel-function series \citep{kinghele1987satellite}; above
that it integrates the same orbit averages exactly over the eccentric
anomaly by 257-node quadrature, because the truncated series raises the
perigee for \(e\) above about 0.3, which is unphysical. Space weather is
held at its value at \(t_0\), using only days that ended before \(t_0\).

The area-to-mass ratio is the dominant unknown, and it is sampled rather
than assumed. Each object draws 32 Monte-Carlo values from one of three
sources, which is recorded per object: its design value where the
catalogue has one, modelled log-normally with a spread of 0.25 dex; a
pool of design values of its class estimated on the training cohort, for
intact payloads and rocket stages without one; or the NASA Standard
Breakup Model \citep{johnson2001nasa} with a characteristic length drawn
from the training distribution of radar-derived sizes of its class, for
fragments and components. Each object has its own random stream seeded
from its catalogue identifier, so a draw is reproducible per object. The
prior is therefore not free of fitted quantities: its class pools are
estimated on the training cohort, and we say so rather than presenting
it as purely a priori.

\subsection{Learned expert}\label{learned-expert}

The expert performs quantile regression of \(z\) on 53 features in seven
groups: orbit, age at \(t_0\), design properties, space weather, family
evidence \(\varphi_F\), five summaries of the physics prior's
Monte-Carlo distribution, and object class. Candidates are models
published between 2022 and 2026: TabM, a parameter-efficient deep
ensemble of multilayer perceptrons sharing most weights
\citep{gorishniy2025tabm}; and two in-context foundation models,
TabICLv2 \citep{qu2025tabicl, qu2026tabiclv2} and Kumo-Tabular
\citep{nvidia2026kumo}, which condition on the training set at inference
time instead of fitting weights to it. The candidate with the lowest
pinball loss on the tuning split becomes the expert. Gradient-boosted
trees
\citep{chen2016xgboost, ke2017lightgbm, prokhorenkova2018catboost}, a
plain non-crossing quantile multilayer perceptron, and an
accelerated-failure-time model trained on all objects including censored
ones \citep{barnwal2022survival} are reported as reference baselines and
are not eligible for selection. Each tuned learner receives the same
budget of 12 random configurations scored on the tuning split and five
seeds; the two in-context models are used with their default inference
settings and are not tuned.

The deliberate exclusion is the radar cross-section. Whether the
catalogue publishes one depends on how fast the object decayed, so its
\emph{availability} carries the outcome; Section \ref{sec:results}
quantifies this and the paper keeps the covariate only as a labelled
leak demonstration.

\subsection{State estimation for objects already in
orbit}\label{state-estimation-for-objects-already-in-orbit}

The module is fitted at canonical-orbit epochs, but a twin must forecast
objects that have been in orbit for years. Survival since the canonical
orbit selects low area-to-mass ratios, which neither the class pools nor
the learned expert know about. For each uncontrolled object the twin
therefore conditions the area-to-mass distribution on the object's
catalogue history, a coarse two-orbit form of the ballistic-coefficient
estimation used in practice
\citep{gondelach2017ballistic, sang2013estimation, pardini2008impact}:

\begin{enumerate}
\def\labelenumi{\arabic{enumi}.}
\tightlist
\item
  128 prior samples are propagated from the canonical orbit to the as-of
  date with the realised trailing space weather;
\item
  samples that would already have re-entered receive zero weight;
\item
  the remainder are weighted by a Gaussian likelihood of the current
  perigee and apogee, with \(\sigma = 10\ \text{km} + 0.3\,|\Delta|\)
  for an observed change \(\Delta\) since the canonical orbit;
\item
  32 posterior samples are drawn by systematic resampling; if the
  effective sample size falls below four, survival alone is used, and if
  no sample survives, the eight longest survivors are.
\end{enumerate}

The physics prior then forecasts from the current orbit with these
samples. This step is Bayesian rather than learned, and Section
\ref{sec:discussion} returns to that point.

\section{Health module}\label{sec:health}

The health module converts a forecaster's residuals into calibrated
alarms. For channel \(c\), a model fitted on nominal data produces a
non-conformity score \(s_c(t)\), which becomes a conformal \(p\)-value
against that channel's own calibration pool \(\mathcal{C}_c\),

\begin{equation}\protect\phantomsection\label{eq:pval}{p_c(t) = \frac{1 + \#\{s' \in \mathcal{C}_c : s' \geq s_c(t)\}}{|\mathcal{C}_c| + 1},}\end{equation}

the \(p\)-values are aggregated across channels at each step with
Benjamini--Hochberg control \citep{benjamini1995}, and an alarm is
raised causally at the third consecutive step with
\(p \leq \max(10^{-3}, 1/(|\mathcal{C}_c|+1))\). The level, the
persistence and the window were fixed before any detector was scored,
and no test label informs any choice. Calibration scores are computed
out of fold in five contiguous blocks with a 30-step purge on each side,
so no step is scored by a model that has seen it.

The detector compared in this paper is Chronos-2
\citep{ansari2025chronos2}, a pretrained time-series foundation model
used zero-shot: rolling eight-step forecasts from up to 512 past steps,
scored by the absolute residual of the forecast median. It is set
against three classic detectors --- ridge autoregression,
principal-component reconstruction and nearest-neighbour window distance
--- the published telemanom scores \citep{hundman2018}, a random-score
floor and an always-on alarm reference.

Two limitations of this evidence must be stated plainly. First, the Soil
Moisture Active Passive and Mars Science Laboratory channels are real
spacecraft telemetry, but they belong to an Earth-observation satellite
and a planetary rover, not to the catalogue population whose debris risk
the rest of the system forecasts; no public dataset provides telemetry
for the objects in our catalogue cohort. The health module therefore
transfers a \emph{method}, not a fitted model, and its coupling to the
debris modules runs through operational status rather than through
telemetry. Second, the conformal alarm layer of this module was
developed and tested in a separate pre-registered study on the European
Space Agency benchmark \citep{kotowski2026esaadb} by the same authors;
that study found the gain on one mission and a reversal on the other,
and we treat the health module as the least validated of the three
throughout.

\section{Cross-module couplings}\label{sec:couplings}

What separates this system from three models in one wrapper is that the
modules exchange specific variables, and that each exchange is measured.
Table \ref{tbl:couplings} states each coupling, what crosses the
boundary, and what the measurement shows.

\begin{table*}[t]
\centering
\footnotesize
\caption{\textbf{The four couplings and their measured effect.}
Confidence intervals are 2,000 cluster bootstraps over launch families
(C1), clouds (C2) or launches (C3). A positive C1 effect means the
features help; a negative C3 difference means the status-aware policy
lowers the Brier score.}\label{tbl:couplings}
\begin{tabular}{@{}
>{\raggedright\arraybackslash}p{(\linewidth - 8\tabcolsep) * \real{0.2000}}
  >{\raggedright\arraybackslash}p{(\linewidth - 8\tabcolsep) * \real{0.2000}}
  >{\raggedright\arraybackslash}p{(\linewidth - 8\tabcolsep) * \real{0.2000}}
  >{\raggedright\arraybackslash}p{(\linewidth - 8\tabcolsep) * \real{0.2000}}
  >{\raggedright\arraybackslash}p{(\linewidth - 8\tabcolsep) * \real{0.2000}}@{}}
\toprule\noalign{}
\begin{minipage}[b]{\linewidth}\raggedright
Coupling
\end{minipage} & \begin{minipage}[b]{\linewidth}\raggedright
Direction
\end{minipage} & \begin{minipage}[b]{\linewidth}\raggedright
What is measured
\end{minipage} & \begin{minipage}[b]{\linewidth}\raggedright
Measured effect
\end{minipage} & \begin{minipage}[b]{\linewidth}\raggedright
Verdict
\end{minipage} \\
\midrule\noalign{}

\bottomrule\noalign{}

C1 & debris \(\rightarrow\) re-entry & family features removed from the
expert & +0.0056 {[}+0.0015, +0.0101{]} at 2016; +0.0024 {[}-0.0058,
+0.0103{]} at 2011 & yes at 2016, not at 2011 \\
C2 & re-entry \(\rightarrow\) debris & cloud population against
catalogue counts & best mean absolute error 0.095 {[}0.068, 0.123{]} of
a cloud (predcal\_tabicl) & experts not separable \\
C3 & health \(\rightarrow\) re-entry & status-aware against status-blind
Brier score & -0.0054 {[}-0.0120, +0.0006{]} & not significant \\
C4 & health \(\times\) debris & no public paired dataset & not evaluated
& demonstration only \\
\end{tabular}
\end{table*}

\textbf{C1, debris to re-entry.} Motivation: fragments of one breakup
share an origin and therefore a range of area-to-mass ratios, so the
fate of catalogued siblings is evidence about an object's own lifetime.
Formulation: \(\varphi_F(i,t_0)\) enters the feature vector \(x_i\) of
the learned expert. Evaluation: the feature group is removed and the
expert refitted with identical settings and seeds, and the paired
difference in test pinball loss is bootstrapped over launch families.
Result: the features help at origin 2016, +0.0056 (+0.0015 to +0.0101),
and the interval includes zero at origin 2011. Limitation: siblings are
only informative once some of them have decayed, which is why the
earlier origin gains less.

\textbf{C2, re-entry to debris.} Motivation: the debris environment is
the sum of individual survival, so a lifetime model should reproduce the
observed decline of a breakup cloud. Formulation: Eq. \ref{eq:cloud}.
Evaluation: predicted against catalogued fragment counts for every cloud
with at least 30 fragments in a test cohort, at one to five years after
its first catalogued fragment, with a bootstrap resampling whole clouds.
Result: the best reading reaches a mean absolute error of about a tenth
of a cloud's size, but with 27 clouds the experts cannot be separated.
Limitation: the number of clouds, not the number of fragments, sets the
precision of this test.

\textbf{C3, health to re-entry.} Motivation: an operator that still
controls a satellite determines its orbit, so drag physics does not
predict its decay. Formulation: the status in \(h_i\) selects a regime
\(\rho_i\); in the operator-controlled regime the drag forecast is
withheld and the previous-window payload base rate is used instead.
Evaluation: an as-of-date test on 13,794 payloads below 600 km with 182
decays in the following 72 days, comparing status-aware and status-blind
policies that use only information available at the as-of date, with
intervals resampling launches. Result: drag physics over-predicts for
operational payloads (mean predicted 0.039 against an observed rate of
0.012) and under-predicts for non-operational ones (0.076 against
0.106). The status-aware policy lowers the Brier score, but the interval
includes zero. Limitation: with 182 decays in one window the test shows
the direction of the effect and not its size.

\textbf{C4, health and debris to response.} Motivation: a high-risk
conjunction for a spacecraft that may be unable to manoeuvre warrants
escalation. Formulation: an alarm state or a non-operational status
raises the response level of a high-risk event. Evaluation: none is
possible, because no public dataset pairs conjunction messages with
telemetry or health status for the same spacecraft. This coupling is
implemented and reported as a demonstration, and it supports no
performance claim anywhere in this paper.

\section{Datasets and preprocessing}\label{sec:data}

All data are public and real; provenance, checksums and licences are
recorded in a manifest released with the code. Table \ref{tbl:datasets}
lists what is used.

\begin{table*}[t]
\centering
\footnotesize
\caption{\textbf{Datasets.} The class column follows the project's own
data audit, which separates real sources from derived and simulated
ones. Simulated sets held on disk, such as the Kelvins debris-origin
challenge data, are not used anywhere in this
paper.}\label{tbl:datasets}
\begin{tabular}{@{}
>{\raggedright\arraybackslash}p{(\linewidth - 8\tabcolsep) * \real{0.2000}}
  >{\raggedright\arraybackslash}p{(\linewidth - 8\tabcolsep) * \real{0.2000}}
  >{\raggedright\arraybackslash}p{(\linewidth - 8\tabcolsep) * \real{0.2000}}
  >{\raggedright\arraybackslash}p{(\linewidth - 8\tabcolsep) * \real{0.2000}}
  >{\raggedright\arraybackslash}p{(\linewidth - 8\tabcolsep) * \real{0.2000}}@{}}
\toprule\noalign{}
\begin{minipage}[b]{\linewidth}\raggedright
Role
\end{minipage} & \begin{minipage}[b]{\linewidth}\raggedright
Source
\end{minipage} & \begin{minipage}[b]{\linewidth}\raggedright
Class
\end{minipage} & \begin{minipage}[b]{\linewidth}\raggedright
Extent used
\end{minipage} & \begin{minipage}[b]{\linewidth}\raggedright
Licence
\end{minipage} \\
\midrule\noalign{}

\bottomrule\noalign{}

Catalogue (re-entry, debris) & General Catalog of Artificial Space
Objects, release 1.8.8; CelesTrak catalogue and space weather & real &
64,516 objects below 2000 km perigee; 32,970 natural decays & Creative
Commons Attribution 4.0 / CelesTrak policy \\
Conjunctions (debris) & European Space Agency Kelvins
collision-avoidance challenge & real, anonymised & 8,293 training and
2,167 official test events & Creative Commons Attribution 4.0 \\
Telemetry (health) & National Aeronautics and Space Administration Soil
Moisture Active Passive and Mars Science Laboratory & real & 81
channels, 68 and 36 labelled sequences & as stated by the authors \\
Thermospheric density & NRLMSIS 2.1 via pymsis, tabulated & derived &
density table used by the physics prior & open source \\
\end{tabular}
\end{table*}

\textbf{Catalogue.} The General Catalog of Artificial Space Objects
\citep{mcdowell2020gcat, mcdowell2022elements} gives, for each object,
one \emph{canonical orbit} representative of early operations with its
epoch, design mass and dimensions, the parent object for fragments, and
the end of the current phase with a status code that distinguishes
natural decay from controlled de-orbit, landing or docking. The
CelesTrak catalogue \citep{celestrak2026satcat} adds operational status,
and the CelesTrak space-weather file provides the solar and geomagnetic
indices; only observed rows for days that ended before the prediction
time are read. The cohort holds 64,516 objects whose canonical orbit has
perigee below 2000 km, of which 32,970 decayed naturally. Figure
\ref{fig:data-reentry} shows the cohort's orbits, its distribution in
time with the cohorts of both forecast origins, the observed lifetimes
with their censoring, and the space-weather record.

\begin{figure*}
\centering
\includegraphics[width=1\linewidth,height=\textheight,keepaspectratio,alt={The catalogue dataset. (a) Canonical orbits by object class; the classes overlap heavily at these densities, and the largest is drawn first. (b) Objects per year of canonical-orbit epoch, with the training, tuning and test cohorts of both origins. (c) Observed time to natural decay, and to censoring for objects still in orbit or whose phase ended otherwise. (d) Solar flux; its 81-day mean is the model input, and the bars mark the two test periods.}]{artwork/Figure_1.png}
\caption{\textbf{The catalogue dataset.} (a) Canonical orbits by object
class; the classes overlap heavily at these densities, and the largest
is drawn first. (b) Objects per year of canonical-orbit epoch, with the
training, tuning and test cohorts of both origins. (c) Observed time to
natural decay, and to censoring for objects still in orbit or whose
phase ended otherwise. (d) Solar flux; its 81-day mean is the model
input, and the bars mark the two test periods.}\label{fig:data-reentry}
\end{figure*}

\textbf{Conjunctions.} The Kelvins collision-avoidance challenge data
\citep{uriot2021cacdataset, esa2019kelvinscac} contain anonymised
conjunction data messages received by the European Space Agency between
2015 and 2019. Following the challenge, inputs are the messages released
at least two days before closest approach and the target is the risk of
the final message. Training events are filtered to mirror the test
inputs, leaving 8,293 training events, of which 66 are high risk, and
the official 2,167 test events, of which 150 are high risk. Figure
\ref{fig:data-debris} shows the breakup clouds of the catalogue and the
structure of the message data.

\begin{figure*}
\centering
\includegraphics[width=1\linewidth,height=\textheight,keepaspectratio,alt={Debris data. (a) Catalogued fragments in orbit over time for the four largest breakup clouds of the cohort. (b) Gabbard diagram of the largest cloud. (c) Risk sequences of 75 training events, 60 drawn at random from the low-risk events and 15 from the high-risk events, so high-risk events are deliberately oversampled. (d) The forecast target: most final risks lie at the floor.}]{artwork/Figure_2.png}
\caption{\textbf{Debris data.} (a) Catalogued fragments in orbit over
time for the four largest breakup clouds of the cohort. (b) Gabbard
diagram of the largest cloud. (c) Risk sequences of 75 training events,
60 drawn at random from the low-risk events and 15 from the high-risk
events, so high-risk events are deliberately oversampled. (d) The
forecast target: most final risks lie at the
floor.}\label{fig:data-debris}
\end{figure*}

\textbf{Telemetry.} The Soil Moisture Active Passive and Mars Science
Laboratory channels \citep{hundman2018} provide 81 channels with 68 and
36 labelled anomaly sequences respectively; the label file lists one
channel twice and the two rows are merged. Figure
\ref{fig:data-telemetry} shows example channels with their labels.

\begin{figure*}
\centering
\includegraphics[width=1\linewidth,height=\textheight,keepaspectratio,alt={Telemetry data. Two channels of each spacecraft drawn at random: training series in grey, test series in black, labelled anomaly sequences shaded. The lower panel shows the European Space Agency benchmark around the first labelled anomaly of its test half, for the target channels it affects.}]{artwork/Figure_3.png}
\caption{\textbf{Telemetry data.} Two channels of each spacecraft drawn
at random: training series in grey, test series in black, labelled
anomaly sequences shaded. The lower panel shows the European Space
Agency benchmark around the first labelled anomaly of its test half, for
the target channels it affects.}\label{fig:data-telemetry}
\end{figure*}

\textbf{What is not available.} Orbital-element histories at catalogue
scale require an account whose user agreement forbids redistribution
\citep{spacetrack2019agreement}, and no legitimate public archive
exists; the European Space Agency's debris database and its re-entry
predictions also need accounts
\citep{klinkrad1991discos, esa2026reentry}. The re-entry module
therefore predicts from a single dated orbit per object rather than from
a history, which is the main structural difference between this study
and operational practice.

\section{Experimental setup and metrics}\label{sec:setup}

\subsection{Temporal protocol}\label{temporal-protocol}

The prediction time \(t_0\) is the canonical-orbit epoch. Under the
protocol \emph{as of origin \(O\)}, the training cohort has
\(t_0 < O - 7\) yr, the tuning cohort \(O - 7 \leq t_0 < O - 5\) yr, and
all their outcomes are censored administratively at \(O\) so that
nothing later than the origin can influence a fitted quantity. The
tuning cohort is split by a hash of the launch family into tune-A, used
for early stopping, the hyperparameter search, learner selection and
gate checks, and tune-B, used for conformal and probability calibration;
splitting by family rather than at random prevents siblings of one
launch from appearing on both sides. The test cohort starts at \(O\) and
ends where the five-year follow-up is complete at the snapshot, or at
the next origin: \(t_0 <\) 2016-01-01 at origin 2011 and \(t_0 <\)
2021-10-27 at origin 2016. The two origins, 2011-01-01 and 2016-01-01,
therefore give two disjoint test periods, which span different phases of
the solar cycle. The primary population is objects that cannot
manoeuvre; payloads are excluded from it because many are
operator-controlled.

For the collision-risk module the official challenge split is used
unchanged. For telemetry the official channel split is used, with
thresholds fixed in advance.

\subsection{Metrics}\label{metrics}

Table \ref{tbl:metrics} lists every metric and what it is for. Pinball
loss over the seven quantiles is primary throughout, because a single
predicted date is not a decision-useful object for either re-entry
planning or conjunction screening.

\begin{table*}[t]
\centering
\footnotesize
\caption{\textbf{Metrics.} Ranking metrics use the raw probabilities,
because isotonic calibration ties scores; calibration metrics use the
mapped ones.}\label{tbl:metrics}
\begin{tabular}{@{}
>{\raggedright\arraybackslash}p{(\linewidth - 4\tabcolsep) * \real{0.3333}}
  >{\raggedright\arraybackslash}p{(\linewidth - 4\tabcolsep) * \real{0.3333}}
  >{\raggedright\arraybackslash}p{(\linewidth - 4\tabcolsep) * \real{0.3333}}@{}}
\toprule\noalign{}
\begin{minipage}[b]{\linewidth}\raggedright
Metric
\end{minipage} & \begin{minipage}[b]{\linewidth}\raggedright
What it measures
\end{minipage} & \begin{minipage}[b]{\linewidth}\raggedright
Where used
\end{minipage} \\
\midrule\noalign{}

\bottomrule\noalign{}

Mean pinball loss & quality of the whole predictive distribution over
seven quantiles & re-entry, conjunction \\
Mean absolute error of the median & point accuracy, \(\log_{10}\) days &
re-entry \\
Empirical coverage, width, Winkler score & honesty and sharpness of the
intervals & re-entry, conjunction \\
Harrell's concordance & ranking over all objects, censored included &
re-entry \\
AUC-PR, AUC-ROC & discrimination at a horizon, on the raw probabilities
& all three modules \\
Brier score, expected calibration error & probability calibration, after
isotonic mapping & re-entry, conjunction \\
Expected over observed count & absolute bias of a probability forecast &
couplings C3, T1 \\
\(L=\mathrm{MSE}_{\mathrm{HR}}/F_2\) & the challenge's official
collision-risk metric & conjunction \\
Corrected event-wise \(F_{0.5}\) & alarm quality, charging alarm time
spent on nominal data & health \\
Kaplan-Meier survival & censoring-aware description of the cohort &
re-entry \\
Cluster bootstrap, 2,000 replicates & intervals resampling launch
families, launches or clouds & all \\
\end{tabular}
\end{table*}

Confidence intervals come from 2,000 cluster-bootstrap replicates,
paired for model differences. The resampling unit is the launch family
for the catalogue cohorts, the launch for the payload-status test and
the cloud for the cloud-population test, because objects sharing an
origin do not provide independent evidence. Learner predictions are
cached under a fingerprint of the data, features and settings, so a
stale result cannot enter a comparison.

\subsection{Deviations and prior use of test
data}\label{deviations-and-prior-use-of-test-data}

The test sets were not untouched when the reported results were
produced, and the full record is in the supplementary material. In
summary: an earlier complete run used a different candidate pool and its
test results were known before the pool was narrowed to models published
2022--2026; two engineering smoke tests printed test scores before the
reported run; every benchmark was scored once before an internal review
and corrected afterwards; the window used for the as-of-date tests of C3
and the state-estimation comparison had been analysed earlier in the
project; and the head-to-head comparisons, the
gate-applied-to-every-learner analysis and the calibrated cloud reading
were added after results had been seen. All intervals are nominal 95 \%
intervals without correction for multiplicity. We report this because it
bounds how confirmatory the comparisons can be, and a fresh test period
would be needed to confirm them.

\section{Results}\label{sec:results}

\subsection{The mixture of experts}\label{the-mixture-of-experts}

Table \ref{tbl:moe} and Figure \ref{fig:moe} give the eight-arm
comparison that tests the paper's central claim.

\begin{table*}[t]
\centering
\footnotesize
\caption{\textbf{Mixture-of-experts arms on both test periods.} Columns
are the mean pinball loss, the mean absolute error of the median in
\(\log_{10}\) days, and the empirical coverage of the 90 \% conformal
interval; the first three numeric columns are origin 2016 and the last
three origin 2011. Lowest pinball loss: 0.0499 at origin 2016 and 0.0527
at origin 2011; the proposed gate attains
neither.}\label{tbl:moe}
\begin{tabular}{@{}
>{\raggedright\arraybackslash}p{(\linewidth - 14\tabcolsep) * \real{0.1250}}
  >{\raggedright\arraybackslash}p{(\linewidth - 14\tabcolsep) * \real{0.1250}}
  >{\raggedleft\arraybackslash}p{(\linewidth - 14\tabcolsep) * \real{0.1250}}
  >{\raggedleft\arraybackslash}p{(\linewidth - 14\tabcolsep) * \real{0.1250}}
  >{\raggedleft\arraybackslash}p{(\linewidth - 14\tabcolsep) * \real{0.1250}}
  >{\raggedleft\arraybackslash}p{(\linewidth - 14\tabcolsep) * \real{0.1250}}
  >{\raggedleft\arraybackslash}p{(\linewidth - 14\tabcolsep) * \real{0.1250}}
  >{\raggedleft\arraybackslash}p{(\linewidth - 14\tabcolsep) * \real{0.1250}}@{}}
\toprule\noalign{}
\begin{minipage}[b]{\linewidth}\raggedright
Fusion arm
\end{minipage} & \begin{minipage}[b]{\linewidth}\raggedright
Type
\end{minipage} & \begin{minipage}[b]{\linewidth}\raggedleft
Pinball
\end{minipage} & \begin{minipage}[b]{\linewidth}\raggedleft
MAE
\end{minipage} & \begin{minipage}[b]{\linewidth}\raggedleft
Cov.
\end{minipage} & \begin{minipage}[b]{\linewidth}\raggedleft
Pinball
\end{minipage} & \begin{minipage}[b]{\linewidth}\raggedleft
MAE
\end{minipage} & \begin{minipage}[b]{\linewidth}\raggedleft
Cov.
\end{minipage} \\
\midrule\noalign{}

\bottomrule\noalign{}

Physics prior only (\(g=1\)) & single expert & 0.0892 & 0.230 & 0.876 &
0.0527 & 0.148 & 0.974 \\
TabM (2025) & single expert & 0.0593 & 0.192 & 0.877 & 0.0994 & 0.354 &
0.918 \\
TabICLv2 (2026) & single expert & 0.0499 & 0.166 & 0.962 & 0.0596 &
0.160 & 0.919 \\
Kumo-Tabular (2026) & single expert & 0.0508 & 0.162 & 0.974 & 0.0616 &
0.163 & 0.966 \\
XGBoost & baseline & 0.0576 & 0.190 & 0.921 & 0.0648 & 0.231 & 0.894 \\
Equal-weight average (\(g=1/2\)) & fixed fusion & 0.0691 & 0.205 & 0.838
& 0.0527 & 0.148 & 0.970 \\
Learned constant gate & learned fusion & 0.0601 & 0.193 & 0.877 & 0.0608
& 0.162 & 0.960 \\
Reliability gate fitted on tune-A & learned fusion & 0.0695 & 0.202 &
0.832 & 0.0616 & 0.163 & 0.959 \\
Reliability-gated mixture (proposed) & learned fusion & 0.0604 & 0.193 &
0.874 & 0.0583 & 0.157 & 0.959 \\
\end{tabular}
\end{table*}

\begin{figure*}
\centering
\includegraphics[width=1\linewidth,height=\textheight,keepaspectratio,alt={All eight fusion arms on both test periods. The dotted line marks the best arm at each origin.}]{artwork_eaai/eaai_moe.png}
\caption{\textbf{All eight fusion arms on both test periods.} The dotted
line marks the best arm at each origin.}\label{fig:moe}
\end{figure*}

The headline comparison is favourable to the learned experts and
unfavourable to the fusion. For 2016--2021 the in-context foundation
models are the most accurate arms, at 0.0499 and 0.0508 against 0.0892
for the physics prior and 0.0576 for the strongest gradient-boosted
baseline, reductions of 44 \% and 13 \% respectively. For 2011--2015 the
ordering changes: the physics prior has the lowest point estimate,
0.0527, and TabM --- the arm selected by the tuning split at the other
origin --- is the worst of all at 0.0994.

The proposed reliability gate reaches 0.0604 and 0.0583. \textbf{It does
not attain the best arm at either origin.} It beats the physics prior
decisively at 2016 and beats its own learned expert at 2011, but an
oracle choosing the single best expert per period would do better in
both. Three further points follow from the table. The equal-weight
average is the second-best arm at 2011 and among the worst at 2016, so a
fixed mixing weight is not a safe default. The learned constant gate,
which is the same model without reliability features, is close to the
full gate at 2016 and worse at 2011, so the reliability descriptors
contribute little at one origin and something at the other. And fitting
the gate on the tuning split rather than on out-of-fold training
predictions is clearly worse at both origins, which shows the fitting
protocol matters more than the parameterisation.

Table \ref{tbl:reentry-pairs} gives the paired comparisons with
intervals.

\begin{table*}[t]
\centering
\footnotesize
\caption{\textbf{Head-to-head differences in test pinball loss} with
paired launch-family cluster bootstraps. A negative difference favours
the first model. These comparisons were added after test results had
been seen and are nominal.}\label{tbl:reentry-pairs}
\begin{tabular}{@{}
>{\raggedright\arraybackslash}p{(\linewidth - 8\tabcolsep) * \real{0.2000}}
  >{\raggedright\arraybackslash}p{(\linewidth - 8\tabcolsep) * \real{0.2000}}
  >{\raggedleft\arraybackslash}p{(\linewidth - 8\tabcolsep) * \real{0.2000}}
  >{\raggedleft\arraybackslash}p{(\linewidth - 8\tabcolsep) * \real{0.2000}}
  >{\centering\arraybackslash}p{(\linewidth - 8\tabcolsep) * \real{0.2000}}@{}}
\toprule\noalign{}
\begin{minipage}[b]{\linewidth}\raggedright
Origin
\end{minipage} & \begin{minipage}[b]{\linewidth}\raggedright
Comparison
\end{minipage} & \begin{minipage}[b]{\linewidth}\raggedleft
\(\Delta\) pinball
\end{minipage} & \begin{minipage}[b]{\linewidth}\raggedleft
95 \% CI
\end{minipage} & \begin{minipage}[b]{\linewidth}\centering
Excludes 0
\end{minipage} \\
\midrule\noalign{}

\bottomrule\noalign{}

2016 & Physics prior only \(-\) TabICLv2 (2026) & +0.0393 & {[}+0.0217,
+0.0622{]} & yes \\
2016 & Physics prior only \(-\) Kumo-Tabular (2026) & +0.0385 &
{[}+0.0203, +0.0613{]} & yes \\
2016 & Physics prior only \(-\) TabM (2025) & +0.0300 & {[}+0.0147,
+0.0493{]} & yes \\
2016 & Physics prior only \(-\) Reliability-gated mixture & +0.0289 &
{[}+0.0148, +0.0469{]} & yes \\
2016 & TabICLv2 (2026) \(-\) Kumo-Tabular (2026) & -0.0008 & {[}-0.0036,
+0.0024{]} & no \\
2016 & TabICLv2 (2026) \(-\) XGBoost & -0.0076 & {[}-0.0121, -0.0036{]}
& yes \\
2016 & Kumo-Tabular (2026) \(-\) XGBoost & -0.0068 & {[}-0.0134,
-0.0012{]} & yes \\
2016 & TabICLv2 (2026) \(-\) Reliability-gated mixture & -0.0104 &
{[}-0.0179, -0.0043{]} & yes \\
2016 & Kumo-Tabular (2026) \(-\) Reliability-gated mixture & -0.0096 &
{[}-0.0181, -0.0023{]} & yes \\
2016 & TabICLv2 (2026) \(-\) TabM (2025) & -0.0093 & {[}-0.0162,
-0.0036{]} & yes \\
2016 & Kumo-Tabular (2026) \(-\) TabM (2025) & -0.0085 & {[}-0.0167,
-0.0017{]} & yes \\
2011 & Physics prior only \(-\) TabICLv2 (2026) & -0.0069 & {[}-0.0217,
+0.0122{]} & no \\
2011 & Physics prior only \(-\) Kumo-Tabular (2026) & -0.0089 &
{[}-0.0253, +0.0133{]} & no \\
2011 & Physics prior only \(-\) TabM (2025) & -0.0468 & {[}-0.1078,
+0.0085{]} & no \\
2011 & Physics prior only \(-\) Reliability-gated mixture & -0.0056 &
{[}-0.0199, +0.0139{]} & no \\
2011 & TabICLv2 (2026) \(-\) Kumo-Tabular (2026) & -0.0020 & {[}-0.0058,
+0.0024{]} & no \\
2011 & TabICLv2 (2026) \(-\) XGBoost & -0.0053 & {[}-0.0124, +0.0001{]}
& no \\
2011 & Kumo-Tabular (2026) \(-\) XGBoost & -0.0032 & {[}-0.0126,
+0.0035{]} & no \\
2011 & TabICLv2 (2026) \(-\) Reliability-gated mixture & +0.0013 &
{[}-0.0016, +0.0044{]} & no \\
2011 & Kumo-Tabular (2026) \(-\) Reliability-gated mixture & +0.0033 &
{[}+0.0007, +0.0055{]} & yes \\
2011 & TabICLv2 (2026) \(-\) TabM (2025) & -0.0399 & {[}-0.0881,
+0.0006{]} & no \\
2011 & Kumo-Tabular (2026) \(-\) TabM (2025) & -0.0378 & {[}-0.0862,
+0.0009{]} & no \\
\end{tabular}
\end{table*}

At 2016 both foundation models are significantly better than the physics
prior, the gated mixture, TabM and XGBoost. At 2011 the physics prior
cannot be separated from either foundation model or from the gated
mixture; the single significant difference among them is that
Kumo-Tabular is \emph{worse} than the gated mixture that fuses it,
+0.0033 (+0.0007 to +0.0055). Figure \ref{fig:gate} shows what the gate
actually does: it places a mean weight of 0.11 and 0.12 on the physics
prior, more when the area-to-mass ratio comes from the breakup model
than when a design value exists, which is the behaviour the descriptors
were designed to produce.

\begin{figure*}
\centering
\includegraphics[width=1\linewidth,height=\textheight,keepaspectratio,alt={Gate behaviour. (a) Mean weight placed on the physics prior, by the provenance of the area-to-mass ratio. (b) The fitted reliability coefficients; a positive coefficient shifts weight towards the physics prior.}]{artwork_eaai/eaai_gate.png}
\caption{\textbf{Gate behaviour.} (a) Mean weight placed on the physics
prior, by the provenance of the area-to-mass ratio. (b) The fitted
reliability coefficients; a positive coefficient shifts weight towards
the physics prior.}\label{fig:gate}
\end{figure*}

\subsection{Model selection under temporal
shift}\label{model-selection-under-temporal-shift}

Selection on a tuning cohort five to seven years older than the test
cohort did not identify the arm that generalised best. At 2016 TabM led
the tuning split by a wide margin and was significantly worse than both
foundation models on test. At 2011 the small tuning cohort (248 objects)
selected Kumo-Tabular, which is tied with TabICLv2 on test. A random
split would have concealed this; only the forward-in-time protocol
exposes it.

\subsection{What the module outputs look
like}\label{what-the-module-outputs-look-like}

Figure \ref{fig:reentry-outputs} shows the predictive distributions for
individual objects and for the cohort.

\begin{figure*}
\centering
\includegraphics[width=1\linewidth,height=\textheight,keepaspectratio,alt={Re-entry module outputs, origin 2016. (a) Forty test objects drawn at random among those with a known five-year outcome: median and 90 \% conformal interval of the physics prior and of the gated mixture, against the observed outcome. (b--d) Predicted median against observed outcome for every such test object.}]{artwork/Figure_6.png}
\caption{\textbf{Re-entry module outputs, origin 2016.} (a) Forty test
objects drawn at random among those with a known five-year outcome:
median and 90 \% conformal interval of the physics prior and of the
gated mixture, against the observed outcome. (b--d) Predicted median
against observed outcome for every such test
object.}\label{fig:reentry-outputs}
\end{figure*}

\subsection{Collision risk}\label{collision-risk-1}

\begin{table*}[t]
\centering
\footnotesize
\caption{\textbf{Collision risk on the official test set.} \(L\) is the
challenge metric, infinite when \(F_2 = 0\) because every median falls
below the high-risk threshold and the metric divides by
zero.}\label{tbl:cac}
\begin{tabular}{@{}
>{\raggedright\arraybackslash}p{(\linewidth - 14\tabcolsep) * \real{0.1250}}
  >{\raggedleft\arraybackslash}p{(\linewidth - 14\tabcolsep) * \real{0.1250}}
  >{\raggedleft\arraybackslash}p{(\linewidth - 14\tabcolsep) * \real{0.1250}}
  >{\raggedleft\arraybackslash}p{(\linewidth - 14\tabcolsep) * \real{0.1250}}
  >{\raggedleft\arraybackslash}p{(\linewidth - 14\tabcolsep) * \real{0.1250}}
  >{\raggedleft\arraybackslash}p{(\linewidth - 14\tabcolsep) * \real{0.1250}}
  >{\raggedleft\arraybackslash}p{(\linewidth - 14\tabcolsep) * \real{0.1250}}
  >{\raggedleft\arraybackslash}p{(\linewidth - 14\tabcolsep) * \real{0.1250}}@{}}
\toprule\noalign{}
\begin{minipage}[b]{\linewidth}\raggedright
Model
\end{minipage} & \begin{minipage}[b]{\linewidth}\raggedleft
Pinball
\end{minipage} & \begin{minipage}[b]{\linewidth}\raggedleft
AUC-PR
\end{minipage} & \begin{minipage}[b]{\linewidth}\raggedleft
AUC-ROC
\end{minipage} & \begin{minipage}[b]{\linewidth}\raggedleft
\(F_2\)
\end{minipage} & \begin{minipage}[b]{\linewidth}\raggedleft
\(L\)
\end{minipage} & \begin{minipage}[b]{\linewidth}\raggedleft
Cov. 90 \%
\end{minipage} & \begin{minipage}[b]{\linewidth}\raggedleft
Cov. high risk
\end{minipage} \\
\midrule\noalign{}

\bottomrule\noalign{}

Persistence + residual climatology & 1.045 & 0.664 & 0.944 & 0.000 &
\(\infty\) & 0.910 & 0.593 \\
XGBoost & 0.855 & 0.606 & 0.945 & 0.000 & \(\infty\) & 0.917 & 0.433 \\
LightGBM & 0.861 & 0.541 & 0.944 & 0.000 & \(\infty\) & 0.905 & 0.253 \\
CatBoost & 0.869 & 0.542 & 0.935 & 0.000 & \(\infty\) & 0.912 & 0.467 \\
Quantile multilayer perceptron & 1.035 & 0.434 & 0.901 & 0.000 &
\(\infty\) & 0.936 & 0.853 \\
TabM (2025) & 0.885 & 0.512 & 0.937 & 0.000 & \(\infty\) & 0.919 &
0.727 \\
TabICLv2 (2026) & 0.692 & 0.665 & 0.961 & 0.057 & 30.06 & 0.918 &
0.327 \\
Kumo-Tabular (2026) & 0.653 & 0.733 & 0.970 & 0.168 & 9.89 & 0.942 &
0.480 \\
Learned constant gate & 0.653 & 0.733 & 0.970 & 0.168 & 9.90 & 0.935 &
0.480 \\
Reliability-gated mixture (proposed) & 0.653 & 0.733 & 0.970 & 0.168 &
9.90 & 0.940 & 0.480 \\
Latest-risk point forecast (challenge baseline) & & & & & 0.69 & & \\
\end{tabular}
\end{table*}

Kumo-Tabular gives the best predictive distribution, 0.653 against 1.045
for persistence, a reduction of 37 \%, and the best discrimination of
high-risk events at AUC-PR 0.733. Table \ref{tbl:cac-pairs} shows the
margin is significant against every comparator.

\begin{table*}[t]
\centering
\footnotesize
\caption{\textbf{Paired bootstrap over test events of the collision-risk
pinball loss.} A negative difference favours
Kumo-Tabular.}\label{tbl:cac-pairs}
\begin{tabular}{@{}
>{\raggedright\arraybackslash}p{(\linewidth - 4\tabcolsep) * \real{0.3333}}
  >{\raggedleft\arraybackslash}p{(\linewidth - 4\tabcolsep) * \real{0.3333}}
  >{\raggedleft\arraybackslash}p{(\linewidth - 4\tabcolsep) * \real{0.3333}}@{}}
\toprule\noalign{}
\begin{minipage}[b]{\linewidth}\raggedright
Comparison
\end{minipage} & \begin{minipage}[b]{\linewidth}\raggedleft
\(\Delta\) pinball
\end{minipage} & \begin{minipage}[b]{\linewidth}\raggedleft
95 \% CI
\end{minipage} \\
\midrule\noalign{}

\bottomrule\noalign{}

Kumo-Tabular \(-\) TabICLv2 & -0.039 & {[}-0.053, -0.025{]} \\
Kumo-Tabular \(-\) XGBoost & -0.202 & {[}-0.233, -0.171{]} \\
Kumo-Tabular \(-\) TabM (2025) & -0.231 & {[}-0.264, -0.199{]} \\
Kumo-Tabular \(-\) Persistence + residual climatology & -0.392 &
{[}-0.451, -0.337{]} \\
\end{tabular}
\end{table*}

On the official point metric, however, no model beats the latest-risk
baseline at 0.694. The reason is visible in Figure
\ref{fig:cac-outputs}: 129 of the 150 high-risk events receive a median
below the threshold, 41 of them between \(-10\) and \(-6\), so \(F_2\)
reaches only 0.17. The conformal interval covers 94.2 \% of all events
but only 48.0 \% of the high-risk ones, because the calibration groups
are coarse exactly where the decisions are made. This reproduces the
challenge's own finding and an independent re-evaluation
\citep{uriot2022spacecraft, vazquez2026latest}.

\begin{figure*}
\centering
\includegraphics[width=1\linewidth,height=\textheight,keepaspectratio,alt={Collision-risk outputs. (a) Final risk against the latest available risk. (b) Final risk against the median forecast. (c--e) Three high-risk test events drawn at random: input messages, predictive quantiles and the final risk.}]{artwork/Figure_11.png}
\caption{\textbf{Collision-risk outputs.} (a) Final risk against the
latest available risk. (b) Final risk against the median forecast.
(c--e) Three high-risk test events drawn at random: input messages,
predictive quantiles and the final risk.}\label{fig:cac-outputs}
\end{figure*}

\subsection{Health}\label{health}

\begin{table*}[t]
\centering
\footnotesize
\caption{\textbf{Telemetry detectors.} AUC-PR is the per-channel mean
over anomalous steps; \(F_{0.5}\) is the corrected event-wise score,
which charges alarm time spent on nominal data. ``Late ref.'' keeps each
detector's amount of alarm time but places it as one block at the end of
the series.}\label{tbl:telemetry}
\begin{tabular}{@{}
>{\raggedright\arraybackslash}p{(\linewidth - 12\tabcolsep) * \real{0.1429}}
  >{\raggedleft\arraybackslash}p{(\linewidth - 12\tabcolsep) * \real{0.1429}}
  >{\raggedleft\arraybackslash}p{(\linewidth - 12\tabcolsep) * \real{0.1429}}
  >{\raggedleft\arraybackslash}p{(\linewidth - 12\tabcolsep) * \real{0.1429}}
  >{\raggedleft\arraybackslash}p{(\linewidth - 12\tabcolsep) * \real{0.1429}}
  >{\raggedleft\arraybackslash}p{(\linewidth - 12\tabcolsep) * \real{0.1429}}
  >{\raggedleft\arraybackslash}p{(\linewidth - 12\tabcolsep) * \real{0.1429}}@{}}
\toprule\noalign{}
\begin{minipage}[b]{\linewidth}\raggedright
Detector
\end{minipage} & \begin{minipage}[b]{\linewidth}\raggedleft
AUC-PR
\end{minipage} & \begin{minipage}[b]{\linewidth}\raggedleft
\(F_{0.5}\)
\end{minipage} & \begin{minipage}[b]{\linewidth}\raggedleft
Late ref.
\end{minipage} & \begin{minipage}[b]{\linewidth}\raggedleft
AUC-PR
\end{minipage} & \begin{minipage}[b]{\linewidth}\raggedleft
\(F_{0.5}\)
\end{minipage} & \begin{minipage}[b]{\linewidth}\raggedleft
Late ref.
\end{minipage} \\
\midrule\noalign{}

\bottomrule\noalign{}

telemanom (published scores) & 0.576 & & & 0.473 & & \\
Chronos-2, zero-shot & 0.343 & 0.170 & 0.197 & 0.367 & 0.035 & 0.224 \\
Ridge autoregression & 0.355 & 0.240 & 0.413 & 0.329 & 0.252 & 0.297 \\
Principal-component reconstruction & 0.322 & 0.249 & 0.306 & 0.328 &
0.212 & 0.314 \\
Nearest-neighbour window & 0.341 & 0.114 & 0.272 & 0.364 & 0.334 &
0.281 \\
Random scores (floor) & 0.129 & 0.000 & 0.000 & 0.141 & 0.000 & 0.000 \\
Always-on alarm (reference) & & 0.000 & & & 0.000 & \\
\end{tabular}
\end{table*}

The zero-shot foundation model does not beat the simple detectors. Its
per-channel AUC-PR is 0.343 and 0.367, level with ridge,
principal-component and nearest-neighbour scoring and far below the
published telemanom scores at 0.576 and 0.473. Two benchmark artefacts
must be stated alongside these numbers. The always-on alarm reaches an
uncorrected event \(F_1\) of 0.548 and 0.537, higher than every
detector, while its corrected \(F_{0.5}\) is zero; and because 93.7 \%
of labelled steps lie in the second half of their series, the late-alarm
reference outscores the detector itself in seven of the eight
detector--spacecraft pairs. Figure \ref{fig:smap-outputs} shows the
scores and alarms on one channel of each spacecraft.

\begin{figure*}
\centering
\includegraphics[width=1\linewidth,height=\textheight,keepaspectratio,alt={Health module outputs on the median channel of each spacecraft by zero-shot AUC-PR. Top: test series and labels. Middle: conformal p-values with the alarm level. Bottom: labels and each detector's causal alarms.}]{artwork/Figure_16.png}
\caption{\textbf{Health module outputs} on the median channel of each
spacecraft by zero-shot AUC-PR. Top: test series and labels. Middle:
conformal \(p\)-values with the alarm level. Bottom: labels and each
detector's causal alarms.}\label{fig:smap-outputs}
\end{figure*}

\section{Ablations and uncertainty evaluation}\label{sec:ablations}

\subsection{Censoring}\label{censoring}

Any evaluation restricted to objects that have already decayed discards
most of the cohort. Table \ref{tbl:km} and Figure \ref{fig:km} give the
Kaplan--Meier description of both test cohorts, which is model-free and
uses the censoring indicator directly.

\begin{table*}[t]
\centering
\footnotesize
\caption{\textbf{Kaplan--Meier description of the test cohorts.}
\(S(\cdot)\) is the estimated survival probability. A median that is
``not reached'' means more than half the group is still in orbit, so no
median lifetime exists for it.}\label{tbl:km}
\begin{tabular}{@{}
>{\raggedright\arraybackslash}p{(\linewidth - 16\tabcolsep) * \real{0.1111}}
  >{\raggedright\arraybackslash}p{(\linewidth - 16\tabcolsep) * \real{0.1111}}
  >{\raggedleft\arraybackslash}p{(\linewidth - 16\tabcolsep) * \real{0.1111}}
  >{\raggedleft\arraybackslash}p{(\linewidth - 16\tabcolsep) * \real{0.1111}}
  >{\raggedleft\arraybackslash}p{(\linewidth - 16\tabcolsep) * \real{0.1111}}
  >{\raggedleft\arraybackslash}p{(\linewidth - 16\tabcolsep) * \real{0.1111}}
  >{\raggedleft\arraybackslash}p{(\linewidth - 16\tabcolsep) * \real{0.1111}}
  >{\raggedleft\arraybackslash}p{(\linewidth - 16\tabcolsep) * \real{0.1111}}
  >{\raggedleft\arraybackslash}p{(\linewidth - 16\tabcolsep) * \real{0.1111}}@{}}
\toprule\noalign{}
\begin{minipage}[b]{\linewidth}\raggedright
Origin
\end{minipage} & \begin{minipage}[b]{\linewidth}\raggedright
Group
\end{minipage} & \begin{minipage}[b]{\linewidth}\raggedleft
\(n\)
\end{minipage} & \begin{minipage}[b]{\linewidth}\raggedleft
Decays
\end{minipage} & \begin{minipage}[b]{\linewidth}\raggedleft
Censored
\end{minipage} & \begin{minipage}[b]{\linewidth}\raggedleft
Median (d)
\end{minipage} & \begin{minipage}[b]{\linewidth}\raggedleft
\(S\)(30 d)
\end{minipage} & \begin{minipage}[b]{\linewidth}\raggedleft
\(S\)(1 yr)
\end{minipage} & \begin{minipage}[b]{\linewidth}\raggedleft
\(S\)(5 yr)
\end{minipage} \\
\midrule\noalign{}

\bottomrule\noalign{}

2016 & all uncontrolled & 2,993 & 1,594 & 46.7\% & 2,780 & 0.895 & 0.768
& 0.558 \\
2016 & fragments & 2,155 & 1,032 & 52.1\% & 3,354 & 0.962 & 0.839 &
0.619 \\
2016 & rocket stages & 385 & 250 & 35.1\% & 1,050 & 0.789 & 0.617 &
0.415 \\
2016 & components & 453 & 312 & 31.1\% & 933 & 0.661 & 0.551 & 0.383 \\
2011 & all uncontrolled & 2,680 & 1,199 & 55.3\% & not reached & 0.921 &
0.845 & 0.744 \\
2011 & fragments & 2,078 & 858 & 58.7\% & not reached & 0.961 & 0.903 &
0.803 \\
2011 & rocket stages & 267 & 157 & 41.2\% & 1,310 & 0.725 & 0.548 &
0.486 \\
2011 & components & 335 & 184 & 45.1\% & 3,233 & 0.823 & 0.708 &
0.571 \\
\end{tabular}
\end{table*}

\begin{figure*}
\centering
\includegraphics[width=1\linewidth,height=\textheight,keepaspectratio,alt={Kaplan--Meier survival of the test cohorts, overall and by object class, with the censored fraction of each group.}]{artwork_eaai/eaai_km.png}
\caption{\textbf{Kaplan--Meier survival of the test cohorts}, overall
and by object class, with the censored fraction of each
group.}\label{fig:km}
\end{figure*}

46.7 \% of the 2016 cohort and 55.3 \% of the 2011 cohort are
right-censored. At origin 2011 the median lifetime of the cohort is
never reached. Rocket stages and components decay substantially faster
than fragments in both periods. This is why the target is capped and
scored only where the five-year outcome is known, why the
accelerated-failure-time baseline is trained on censored observations as
well, and why the comparison in Table \ref{tbl:gap} against studies that
train on decayed objects only is not like-for-like in their favour.

\subsection{Feature groups}\label{feature-groups}

Removing each feature group from the selected expert and refitting with
identical settings and seeds gives the paired effects in Figure
\ref{fig:ablation}. The family features are the coupling C1 result
reported in Table \ref{tbl:couplings}. Design properties help at 2016
(+0.0071, +0.0035 to +0.0121); the physics summaries help the selected
expert at 2011 (+0.0022) but hurt at 2016 (-0.0050), which is consistent
with the physics prior's own error growing in the rising phase of the
solar cycle. All intervals are nominal and uncorrected for the eight
comparisons.

\begin{figure*}
\centering
\includegraphics[width=0.78\linewidth,height=\textheight,keepaspectratio,alt={Feature-group ablations of the selected expert. In the first four rows a group is removed, so a positive change means the group helps; the last row instead adds the leaking covariate, where a positive change means it makes the model worse.}]{artwork/Figure_9.png}
\caption{\textbf{Feature-group ablations of the selected expert.} In the
first four rows a group is removed, so a positive change means the group
helps; the last row instead adds the leaking covariate, where a positive
change means it makes the model worse.}\label{fig:ablation}
\end{figure*}

\subsection{Leakage through covariate
availability}\label{leakage-through-covariate-availability}

Whether the catalogue publishes a radar cross-section depends on how
quickly the object decayed. In the origin-2016 training cohort 75.9 \%
of objects without a value decayed within 30 days of their
canonical-orbit epoch, against 17.2 \% of those with one; adjusted for
orbit, class and decade, and with standard errors clustered on launch
families, the odds ratio is 3.59 (2.09 to 6.17).

The consequence for a modeller is instructive. Adding the covariate
changes the test pinball loss by amounts whose intervals include zero at
both origins, so point accuracy would not reveal the problem. What it
does break is calibration: at origin 2016, where the missingness pattern
differs from training, the 90 \% coverage of the selected expert falls
from 87.7 \% to 76.8 \%. A covariate that looks like a size measurement
is in fact partly an outcome indicator, and the symptom appears in the
interval rather than in the mean. We therefore exclude it, and release
the availability test with the code.

\subsection{Uncertainty
quantification}\label{uncertainty-quantification}

\begin{table*}[t]
\centering
\footnotesize
\caption{\textbf{Interval rules on both test periods.} ``Cov. all'' is
over every test object including those surviving past the five-year cap;
``Cov. decayed'' is restricted to objects that decayed within five
years; width is in \(\log_{10}\) days. The first three numeric columns
are origin 2016, the last three origin
2011.}\label{tbl:coverage}
\begin{tabular}{@{}
>{\raggedright\arraybackslash}p{(\linewidth - 12\tabcolsep) * \real{0.1429}}
  >{\raggedleft\arraybackslash}p{(\linewidth - 12\tabcolsep) * \real{0.1429}}
  >{\raggedleft\arraybackslash}p{(\linewidth - 12\tabcolsep) * \real{0.1429}}
  >{\raggedleft\arraybackslash}p{(\linewidth - 12\tabcolsep) * \real{0.1429}}
  >{\raggedleft\arraybackslash}p{(\linewidth - 12\tabcolsep) * \real{0.1429}}
  >{\raggedleft\arraybackslash}p{(\linewidth - 12\tabcolsep) * \real{0.1429}}
  >{\raggedleft\arraybackslash}p{(\linewidth - 12\tabcolsep) * \real{0.1429}}@{}}
\toprule\noalign{}
\begin{minipage}[b]{\linewidth}\raggedright
Interval rule
\end{minipage} & \begin{minipage}[b]{\linewidth}\raggedleft
Cov. all
\end{minipage} & \begin{minipage}[b]{\linewidth}\raggedleft
Cov. decayed
\end{minipage} & \begin{minipage}[b]{\linewidth}\raggedleft
Width
\end{minipage} & \begin{minipage}[b]{\linewidth}\raggedleft
Cov. all
\end{minipage} & \begin{minipage}[b]{\linewidth}\raggedleft
Cov. decayed
\end{minipage} & \begin{minipage}[b]{\linewidth}\raggedleft
Width
\end{minipage} \\
\midrule\noalign{}

\bottomrule\noalign{}

Physics median ±20 \% (operational rule) & 0.602 & 0.122 & 0.18 & 0.747
& 0.171 & 0.17 \\
Physics median ±30 \% & 0.638 & 0.204 & 0.27 & 0.766 & 0.234 & 0.27 \\
Physics prior, conformal 80 \% & 0.738 & 0.419 & 0.35 & 0.944 & 0.805 &
0.65 \\
Physics prior, conformal 90 \% & 0.876 & 0.726 & 0.67 & 0.974 & 0.906 &
1.03 \\
Gated median ±20 \% & 0.529 & 0.183 & 0.18 & 0.651 & 0.201 & 0.18 \\
Gated median ±30 \% & 0.599 & 0.280 & 0.27 & 0.694 & 0.283 & 0.27 \\
Gated mixture, conformal 80 \% & 0.717 & 0.543 & 0.43 & 0.923 & 0.776 &
0.88 \\
Gated mixture, conformal 90 \% & 0.874 & 0.835 & 0.67 & 0.959 & 0.884 &
1.21 \\
\end{tabular}
\end{table*}

Against the operational practice it would replace, the calibrated
interval is a clear improvement on the objects that matter. The ±20 \%
residual-lifetime window around the physics median --- the rule assessed
at roughly 90 \% coverage in the final days before re-entry
\citep{pardini2018assessing} --- covers 12.2 \% and 17.1 \% of objects
that decayed within five years when it is applied at canonical-orbit
epochs, while the conformal 90 \% interval covers 72.6 \% and 90.6 \%.
The comparison is not a criticism of the operational rule, which is
designed for a different regime with continuous tracking; it quantifies
how far that regime is from forecasting years ahead from one catalogue
orbit. Figure \ref{fig:coverage} places every rule side by side.

\begin{figure*}
\centering
\includegraphics[width=0.92\linewidth,height=\textheight,keepaspectratio,alt={Interval coverage over all test objects, including those still in orbit after five years.}]{artwork/Figure_8.png}
\caption{\textbf{Interval coverage} over all test objects, including
those still in orbit after five years.}\label{fig:coverage}
\end{figure*}

\textbf{Conditional coverage.} Mondrian grouping is the mechanism
intended to make coverage conditional, and Table \ref{tbl:coverage}
shows it reaching the nominal level marginally at origin 2011 and
falling short at origin 2016 for the arms with the narrowest intervals.

\textbf{Temporal calibration.} Conformal prediction guarantees
\emph{marginal} coverage under exchangeability, and a forward-in-time
test period is precisely where exchangeability is doubtful. Table
\ref{tbl:temporal} and Figure \ref{fig:temporal} break the 90 \%
coverage down by calendar year of the prediction time.

\begin{table*}[t]
\centering
\footnotesize
\caption{\textbf{Conditional coverage over time.} Years with fewer than
25 scored objects are omitted.}\label{tbl:temporal}
\begin{tabular}{@{}
>{\raggedright\arraybackslash}p{(\linewidth - 6\tabcolsep) * \real{0.2500}}
  >{\raggedright\arraybackslash}p{(\linewidth - 6\tabcolsep) * \real{0.2500}}
  >{\raggedright\arraybackslash}p{(\linewidth - 6\tabcolsep) * \real{0.2500}}
  >{\raggedleft\arraybackslash}p{(\linewidth - 6\tabcolsep) * \real{0.2500}}@{}}
\toprule\noalign{}
\begin{minipage}[b]{\linewidth}\raggedright
Origin
\end{minipage} & \begin{minipage}[b]{\linewidth}\raggedright
Model
\end{minipage} & \begin{minipage}[b]{\linewidth}\raggedright
90 \% coverage, first to last test year
\end{minipage} & \begin{minipage}[b]{\linewidth}\raggedleft
First \(-\) last
\end{minipage} \\
\midrule\noalign{}

\bottomrule\noalign{}

2016 & Reliability-gated mixture & 0.902 · 0.963 · 0.848 · 0.879 · 0.812
· 0.847 & +0.055 \\
2016 & Physics prior only & 0.929 · 0.974 · 0.908 · 0.885 · 0.836 ·
0.732 & +0.197 \\
2016 & TabM (2025) & 0.897 · 0.963 · 0.823 · 0.868 · 0.817 · 0.886 &
+0.011 \\
2011 & Reliability-gated mixture & 0.978 · 0.974 · 0.929 · 0.930 · 0.957
& +0.021 \\
2011 & Kumo-Tabular (2026) & 0.975 · 0.982 · 0.929 · 0.943 · 0.968 &
+0.007 \\
2011 & Physics prior only & 0.980 · 0.980 · 0.965 · 0.957 · 0.978 &
+0.002 \\
\end{tabular}
\end{table*}

\begin{figure*}
\centering
\includegraphics[width=1\linewidth,height=\textheight,keepaspectratio,alt={Coverage by calendar year. The dashed line is the nominal level. Coverage drifts downward across the 2016 test period and is stable across the 2011 period.}]{artwork_eaai/eaai_temporal.png}
\caption{\textbf{Coverage by calendar year.} The dashed line is the
nominal level. Coverage drifts downward across the 2016 test period and
is stable across the 2011 period.}\label{fig:temporal}
\end{figure*}

At origin 2016 the physics prior's coverage falls from 0.929 in the
first test year to 0.732 in the last, a drop of 0.197, while at origin
2011 it stays between 0.978 and 0.980. The marginal guarantee is kept;
what degrades is coverage conditional on time, as the rising phase of
solar cycle 25 moves the covariate distribution away from the
calibration cohort. To our knowledge this is the first demonstration of
the effect on an operational space-surveillance catalogue, and it is a
caution for any deployment that calibrates once and then runs for years.

\subsection{State estimation for objects in
orbit}\label{state-estimation-for-objects-in-orbit}

The lifetime pipeline is fitted at canonical-orbit epochs; applied to
objects that have survived for years it over-predicts re-entry by
9.5-fold over a 72-day horizon (6.9 to 13.0), because survival selects
low area-to-mass ratios that the fitted model cannot know about. Class
priors reduce this to 3.64-fold and conditioning on the catalogue
history to 1.37-fold (1.20 to 1.63), still significantly above one. The
three variants rank decaying objects almost identically, with AUC-ROC
within 0.01 of each other; what differs is calibration, which is the
quantity the state-estimation step targets. Figure \ref{fig:t1} shows
the reliability and the counts.

\begin{figure*}
\centering
\includegraphics[width=0.92\linewidth,height=\textheight,keepaspectratio,alt={As-of-date test of the three variants for objects in orbit on 2026-08-16, scored against decays to 2026-10-27. (a) Reliability in six fixed probability bins, marker area growing with the number of objects. (b) Expected against observed re-entries.}]{artwork/Figure_14.png}
\caption{\textbf{As-of-date test of the three variants} for objects in
orbit on 2026-08-16, scored against decays to 2026-10-27. (a)
Reliability in six fixed probability bins, marker area growing with the
number of objects. (b) Expected against observed
re-entries.}\label{fig:t1}
\end{figure*}

This window had been analysed earlier in the project, so this comparison
is an as-of-date test rather than a blind one; a dated, hashed forecast
released with the code is the blind test, and a later catalogue snapshot
will score it.

\section{Discussion, limitations and threats to
validity}\label{sec:discussion}

\subsection{What the evidence
supports}\label{what-the-evidence-supports}

Three claims survive the measurements. First, \textbf{in-context tabular
foundation models transfer to an operational space catalogue}: at origin
2016 they reduce pinball loss by 13 \% against the best gradient-boosted
baseline and by 44 \% against the physics prior, with intervals
excluding zero, and the same family of models gives significantly the
best collision-risk distribution. This is a positive result for
in-context learning outside its usual benchmark setting. Second,
\textbf{distribution-free calibration is what makes the output usable}:
the conformal interval covers a far larger share of decayed objects than
the operational window applied at the same horizon, and the Mondrian
grouping is what keeps coverage near nominal at the earlier origin.
Third, \textbf{the state estimate matters more than the learner} for
objects already in orbit, where a Bayesian conditioning step removes
most of a ninefold bias that no choice of learner addressed.

\subsection{What the evidence does not
support}\label{what-the-evidence-does-not-support}

The reliability gate, which is the architectural novelty, \textbf{does
not beat the best single expert at either origin}. Table \ref{tbl:moe}
makes this explicit, and we draw three conclusions rather than one. The
gate is fitted on out-of-fold predictions of the training period, so it
can learn where the physics prior is unreliable but cannot anticipate
that a \emph{learner} will degrade on a future cohort; its descriptors
are all properties of the prior. The equal-weight and constant-gate arms
show that the benefit of learning the weight at all is small and
inconsistent across origins. And the tuning-split variant shows that the
fitting protocol dominates the parameterisation. A gate that also saw
descriptors of covariate shift is the obvious next step, and it would
have to be validated on a third test period rather than on these two.

The honest reading is therefore that reliability gating is a safety
mechanism rather than an accuracy mechanism: it prevented the
catastrophic arm at each origin (TabM at 2011, the physics prior at
2016) without matching the best one. For an operator who must fix an
architecture before knowing which regime the next five years bring, that
property has value; for a benchmark table it is a loss.

Two further negative results are reported. Zero-shot time-series
forecasting did not beat ridge autoregression on spacecraft telemetry,
and on the official collision-risk metric no model beat a baseline that
simply repeats the latest risk.

\subsection{Threats to validity}\label{threats-to-validity}

\textbf{Test-set reuse.} The test sets were scored more than once during
development, the candidate pool was narrowed after an earlier run's test
results were known, and several analyses were added after results had
been seen. Section \ref{sec:setup} states this and the supplement
records it in full. This bounds the comparisons' confirmatory value; the
significance statements should be read as nominal.

\textbf{Two periods only.} Every statement about generalisation rests on
two test periods from one catalogue. The direction of the solar-cycle
effect is measured, not assumed, but a third period would be needed to
establish that the foundation models keep their advantage.

\textbf{Single orbit per object.} The re-entry module predicts from one
dated orbit because orbital histories at catalogue scale cannot be
redistributed. Operational tools that fit a ballistic coefficient to
weeks of tracking data use far more information near re-entry, and our
months-to-years forecasts are not comparable with their final-days
accuracy.

\textbf{Priors estimated from catalogue values.} The physics prior's
class pools are built from objects that received a radar-derived size,
which biases them towards longer-lived objects; the design-data ablation
bounds the influence of this.

\textbf{Telemetry domain.} The telemetry benchmarks are real but belong
to spacecraft outside the catalogue population, so the health module
transfers a method and not a fitted model, and its coupling into the
debris modules runs through operational status rather than telemetry.
The benchmark's own metric artefacts, documented in Section
\ref{sec:results}, mean its absolute scores should not be read as
operational detection rates.

\textbf{Integrator resolution.} A step of the decay integrator may not
exceed ten days, so for an object whose perigee is already within one
kilometre of the re-entry threshold the forecast is reported at the end
of that step; two objects in the released forecast have a degenerate
interval for this reason. They are within days of re-entry regardless,
and no reported comparison depends on them.

\textbf{Reproducibility.} Seeded training on the graphics processor is
not bit-reproducible, while on the central processor it is. Repeating
the telemetry run moved per-channel AUC-PR by at most
\(2 \times 10^{-5}\) and changed no reported value at three decimals;
two runs of the lifetime variant in the as-of-date test gave ratios of
8.9 and 9.5.

\section{Conclusions and future work}\label{sec:conclusion}

This paper presented a shared-state digital twin for orbital debris
analysis in which a reliability-gated mixture of experts, calibrated by
conformal prediction and isotonic regression, forecasts residual orbital
lifetime and collision risk, and in which spacecraft health enters
through operational status. It was evaluated on 64,516 catalogued
objects across two forward-in-time test periods, an official
collision-avoidance benchmark and public spacecraft telemetry, with
cluster-bootstrap intervals throughout.

The measured findings are:

\begin{enumerate}
\def\labelenumi{\arabic{enumi}.}
\tightlist
\item
  In-context tabular foundation models are the most accurate lifetime
  arms in the later test period, by 13 \% against the best
  gradient-boosted baseline and 44 \% against the physics prior, and
  they give significantly the best collision-risk distribution, 37 \%
  better than persistence. In the earlier period they are statistically
  tied with the physics prior.
\item
  The reliability gate did not beat the best single expert at either
  origin. It avoided the worst arm in both, and the fitting protocol
  mattered more than the gate's parameterisation.
\item
  Calibrated intervals cover 72.6 \% and 90.6 \% of decayed objects
  where the operational ±20 \% window applied at the same horizon covers
  12.2 \% and 17.1 \%.
\item
  Conformal coverage conditional on time degrades across a six-year
  forward test window, falling by 0.197 at origin 2016 while remaining
  stable at origin 2011.
\item
  Conditioning the area-to-mass ratio on an object's catalogue history
  reduces 72-day over-prediction from 9.5-fold to 1.37-fold.
\item
  Of the four couplings, breakup-family evidence helps the learned
  expert in one of two periods, cloud population is reproduced to about
  a tenth of a cloud without separating the experts, the status coupling
  has the expected direction but no significant gain, and the fourth
  cannot be tested on public data.
\item
  A catalogue covariate can leak the outcome through its availability.
  Adding it leaves point accuracy unchanged within noise and breaks
  interval coverage, so leakage of this kind is invisible to a
  mean-error check.
\end{enumerate}

Future work follows directly from the negative results. The gate should
be given descriptors of covariate shift and validated on a third period.
The state-estimation step, which produces the largest effect in the
paper, is currently Bayesian importance sampling over a physics
propagator; replacing it with an amortised neural posterior trained on
simulated propagations would make the system's strongest component a
learned one, and we are pursuing that. On the collision-risk side, the
gap between the best predictive distribution and the official point
metric is a decision-theoretic problem rather than a modelling one, and
a learned decision rule over the calibrated probability is the natural
next step. Finally, the dated and hashed forecast released with this
paper is the only fully blind test of the system, and scoring it against
a later catalogue snapshot is the cleanest available validation.

\section*{Declarations}\label{declarations}
\addcontentsline{toc}{section}{Declarations}

\textbf{Contribution in artificial intelligence and application in
engineering.} The contribution in artificial intelligence is the
reliability-gated mixture of experts with distribution-free calibration
described in Section \ref{sec:arch}, together with the evaluation of
recent tabular and time-series foundation models under forward-in-time
drift and administrative censoring. The application in engineering is
forecasting the low-Earth-orbit debris environment: residual orbital
lifetime, collision risk and spacecraft health status.

\textbf{Credit authorship contribution statement.} {[}Author 1{]}:
conceptualisation, methodology, software, validation, formal analysis,
investigation, data curation, writing --- original draft, visualisation.
{[}Author 2{]}: supervision, writing --- review and editing. {[}To be
completed by the authors.{]}

\textbf{Declaration of competing interest.} {[}To be completed by the
authors.{]}

\textbf{Funding.} {[}To be completed by the authors.{]}

\textbf{Data availability.} All data used are public. The catalogue is
available under Creative Commons Attribution 4.0
\citep{mcdowell2020gcat}; the CelesTrak catalogue and space-weather
files under CelesTrak's usage policy \citep{celestrak2026satcat}; the
collision-avoidance challenge data under Creative Commons Attribution
4.0 \citep{uriot2021cacdataset}; and the telemetry benchmarks under the
terms stated by their authors \citep{hundman2018, esaadb_dataset}. The
derived cohort table, the density table, all predictions and result
files, and the dated forecast will be deposited with a digital object
identifier before publication {[}identifier to be added{]}. No
restricted orbital-element data are used or redistributed. Pretrained
weights are obtained from their publishers under their own licences.

\textbf{Code availability.} The code, including data fetchers with
checksummed manifests, the three modules, the twin engine and the
scripts that generate every table and figure in this paper from the
result files, is available \citep{cosmos_software}; an archived release
with a digital object identifier will be created before publication.

\textbf{Declaration of generative artificial intelligence in the writing
process.} During the preparation of this work the authors used a
generative artificial-intelligence assistant to help write analysis and
plotting code and to draft and edit text. The authors reviewed and
edited the content and take full responsibility for the published
article. Every figure was rendered by the scripts named above from the
study's result files; no figure was generated by a generative model.
{[}The authors must review this statement and adjust it to their actual
use.{]}

\bibliographystyle{elsarticle-num}
\bibliography{refs}

\end{document}
